\section{How to address the \nn{} problem}
\label{sec:fix}

The fixes outlined here would far easier done at the start of the project, and not after the customer hits the issue.
That said these fixes can still be applied to existing operators.

\subsection{Audit the RBAC}
\label{sec:fix:audit}
This is the first step that should be done, and the audit is not about the permission, but for the resources being accessed.
I am working with the assumption that the RBAC has only the CRUD permissions that is it requires, and not overprivileged.
This audit is to classify the resources that the operator requires to function.
There are three classifications that all resources will fall into.

\begin{itemize}
	\item Native kebernetes resource.
	\item Third party CRDs.
	\item Operator native CRDs.
\end{itemize}

Once the resources have being classified a set of test can be defined to see the impact on the operator.
There is a level of risk associated with the different types of resources.
In the following sections that risk is outlined of the resource classification from above.

\subsubsection{Native kebernetes resource}
The native kebernetes resources are generic to every cluster, and this fact makes them the worst problem kind.
Every other operator on the cluster most likely interacts with the native resources, meaning your operator needs the toughest cache filtering here.
These are resources kebernetes itself will create.
It is possible for \textit{everything} on the cluster to create these resources.

\subsubsection{Third party CRDs}
Third party CRDs are those defined by other operators that your operator may integrate with.
An example here would the CRDs for the cluster's gateway that is possible using the gateway API\cite{noauthor_gateway_nodate}.
These resources maybe used by other operators in the cluster, with the other operators defining their own CRs\footnote{A Custom Resource (CR) is a declarative Kubernetes API object, instantiated from a Custom Resource Definition (CRD), that specifies the desired configuration and state of an operator-managed operand.}.
This is a resource that is possible for your operator to be a \nn{} on, and there should be a strong level of cache filtering carried out here, but the third party CRDs are not as risky as the native kebernetes resources.


\subsubsection{Operator native CRDs}
The operator native CRDs are the simplest of the three resource types.
Most likely in the operator that defines the CRDs, the operator is going to want to act upon every instance of that resources.
This would mean the operator native CRDs would have the least caching filtering, if any at all.


\subsection{Label all resources}
\label{sec:fix:label}
If the operator has any interaction with a resource, add a unique label\cite{noauthor_labels_nodate} to that resource related to the operator.
I would raise a level of caution here.
Depending on how the operator is installed different labels maybe applied during the different install methods.
Having a common label across all the resources makes configuring the cache easier.

One problem with adding labels, then filtering the cache on those labels.
If someone removes the label from the resource, your operator does not get an event for that update to the resource.
In the eyes of the operator the resources doesn't exist, but also it does not know it should recreate the resource.

\subsection{Create test scripts}
In one sense this problem is hard to automate tests for, but that does not mean the test can not be scripted.
To test for the \nn{} all you need to do is add native kebernetes resources and third party resources to a namespace your operator should not care about.
By watching the memory usage of your operator you can see if your operator is affected.
Don't only add a few, add a lot.
This will show the issue much faster.

For example one of the operators I was working on had a watch on \textit{deployments}\cite{noauthor_deployments_nodate}.
A simple script\cite{noauthor_kuadrant-crc-setupcreate-deploymentsfish_nodate}  that would create a deployment with the number of pod replica's at zero, and importantly a large annotations field (about 190kbs in size) would highlight the issue.
When the script creates about 500 deployments, the operator's memory would jump by around 200MiB.
That jump is really easy to see on a graph, and very repeatable.
The repeatable allows validation of the caching configuration once implemented.

A word of warning on creating these scripts, and possibly using them to create a large set of resources.
While it will validate that your operator is not a \nn{}, it doesn't mean other operators in your cluster are not.
It can be very possible to crash the entire cluster when proving your operator is not affected.

\subsection{Configure the controller-runtime cache}
Once the resources that have being found during the audit section (\cref{sec:fix:audit}) the controller-runtime cache should be configured to limit what resources are cached.
This should make use of the labels defined in \cref{sec:fix:label}.

A second configuration that is useful to configure is limiting the client that controller-runtime creates to only get resources with existing watches configured using \textit{ReaderFailOnMissingInformer}\cite{noauthor_cache_nodate}.
This will stop developers accidentally from setting up watches for resources that maybe only accessed in the one of cases.

\subsection{Review Again and Again}
Overtime as the operator is build to do more complex work, the chance of new APIs added to the RBAC is likely.
The review of the RBAC, and controller-runtime cache configuration needs to be done on every update.
With the RBAC rules needing update to allow new APIs this can be used as an indicator of when a review is required.
Workflows could be configured to trigger on the RBAC rules to be changed.
At the time of writing inc I have not seen examples of such a workflow, or a successful review process.

